%% kapa: exam=1 round=24 year=2013 period=1 form=A subject=economics subject_ko=경제원론 seq=17 number=57
%% source: Q-Net 감정평가사 기출문제 (cst00309) artlSeq=5073802 file=extracted/1008/5073802/2013년 제24회 감정평가사 1차시험 기출문제/1교시(A형).hwp
%% conversion: kordoc 4.24.51 -> tools/gen1.py (auto); review=proofread-sonnet; flags=none
\begin{question}{57}
평균총비용곡선이 U자형인 A기업의 단기 비용곡선에 관한 설명으로 옳은 것은?（단, 생산요소가격은 불변이며, 고정생산요소가 존재한다.)
\begin{choices}
  \item 한계비용이 평균총비용보다 클 때, 생산량이 증가함에 따라 평균총비용은 증가한다.
  \item 한계비용이 평균총비용보다 클 때, 생산량이 증가함에 따라 평균고정비용은 증가한다.
  \item 평균총비용곡선과 평균가변비용곡선은 동일한 생산량 수준에서 최저점에 도달한다.
  \item 어떤 생산량의 한계비용은 원점에서 그 생산량에 해당하는 가변비용곡선상의 점을 이은 직선의 기울기이다.
  \item 어떤 생산량의 평균총비용은 그 생산량에 해당하는 총비용곡선상의 점에서 그은 접선의 기울기이다.
\end{choices}

eviewflag{choices 2-5 not rendered in rhwp page; only stem and choice 1 verified}
\end{question}
\answer{1}
